\section{Riproducibilità}\label{sec:rip}

\paragraph{Tipi di prova.} Ogni cifra del documento appartiene a uno di quattro tipi: dimostrazione in un modello dichiarato; calcolo in forma chiusa; simulazione; elaborazione di dati pubblici. Le simulazioni sono di due specie: con modello ideale del sensore (misure di posizione con rumore gaussiano e ombra geometrica, senza immagini di profondità) e su profondità simulata dalle schede tecniche di tre camere a stereo attivo (RealSense D435i e D455, Orbbec Gemini 335), a due livelli di rumore, nominale e limite. I dati pubblici sono traiettorie di pedoni al chiuso con verità di riferimento \cite{thor}, che non contengono profondità, sequenze di profondità registrate \cite{tum-rgbd,depthtrack} e una registrazione di esempio distribuita dal costruttore di una delle camere, usata solo per il confronto dei pixel senza dato (\S\ref{sec:val}). Nessuna prova usa un sensore fisico.

\paragraph{Linguaggi.} Le simulazioni con modello ideale del sensore e il predittore sono in C (compilatore gcc 16.2.1, opzione \texttt{-O2}); la catena di riferimento su profondità simulata, la regola del «libero» della cella e i calcoli in forma chiusa sono in Python con una libreria numerica. I tempi di calcolo sono misurati su un calcolatore portatile con processore AMD Ryzen AI 9 HX 370 e sistema Linux, con un solo filo d'esecuzione; non sono tempi della scheda di bordo, e le misure rifatte a macchina carica non sono stabili (\S\ref{sec:cos}).

\paragraph{Generatore di numeri casuali e semi.} Le simulazioni in C usano un generatore a incremento costante: stato $s\in\mathbb{Z}_{2^{64}}$, a ogni estrazione $s\leftarrow s+G \pmod{2^{64}}$ con $G$ dispari, uscita $\mu(s)$ con $\mu$ funzione di mescolamento biiettiva. Lo stato iniziale del seme $\sigma$ è $s_0(\sigma)=\mu(\sigma G+c)$. La forma senza mescolamento, $s_0(\sigma)=\sigma G+c$, rende il flusso del seme $\sigma+1$ uguale a quello del seme $\sigma$ avanzato di un'estrazione (1\,000 valori uguali su 1\,000; con il mescolamento 0 su 1\,000), e gli intervalli binomiali calcolati su semi contigui risultano troppo stretti di un fattore fra $1{,}3$ e $2{,}5$ (\S\ref{sec:val}). Le cifre di continuità e di copertura sono calcolate con semi inizializzati nella forma corretta; taratura e verifica usano insiemi di semi disgiunti. In nessuna prova la libreria di taratura viene da un generatore di camminate e i passaggi di verifica da un altro.

\paragraph{Numerosità tipiche.} Continuità e copertura con modello ideale del sensore: 23\,402 passaggi nascosti di 6\,000 camminatori distinti; sull'insieme di tutti i lanci, 217\,861 passaggi confermati su 226\,400 lanciati, di al più 70\,400 camminatori distinti; da 608 a 2\,359 casi per cella della tavola di copertura. Regola del «libero» della cella: 10 scene, 300 condizioni, 9\,000 fotogrammi, in due lanci a semi diversi. Profondità simulata dalle schede tecniche: 90 passaggi (3 disposizioni per 30 semi) per le prove di base, 180 passaggi per risoluzione per l'identità al ritorno. Traiettorie pubbliche di pedoni: 219 tratti di 9 persone, con numerosità efficace fra 9 e 219. Regola delle etichette: 20\,000 sequenze, due volte. Durata delle esecuzioni sul calcolatore indicato: circa 4 minuti per la simulazione d'insieme, 2 minuti per la regola del «libero», da 12 a 15 minuti per le prove d'identità.

\paragraph{Come si rifanno i conti principali.} Dal solo testo si ricalcolano:
\begin{enumerate}[label=(\alph*)]
\item l'intervallo di Wilson al $95\,\%$ \cite{wilson1927,brown2001} di una frequenza $k/N$: per $k=826$, $N=879$ si ottiene $[0{,}9220;\,0{,}9536]$; per zero eventi su $N$ il limite superiore è $z^2/(N+z^2)$ con $z^2=3{,}8415$, cioè $6{,}4$ su $10\,000$ per $N=6\,000$;
\item la velocità coperta dalla fascia di guardia, $w/(0{,}133+0{,}155+0{,}067)\unita{s}=0{,}40/0{,}355=1{,}13\unita{m/s}$, e con $0{,}20\unita{m}$ d'errore del bordo $(0{,}40-0{,}20)/0{,}355=0{,}56\unita{m/s}$;
\item la velocità coperta dell'avviso a sensore fermo, $(3-0{,}95)/(0{,}267+0{,}155+T_r)$: $1{,}44\unita{m/s}$ con $T_r=1{,}0\unita{s}$, $2{,}84$ con $0{,}3$, $1{,}07$ con $1{,}5$, $0{,}92$ con $1{,}8$ (\S\ref{sec:lim-portatore});
\item la somma di latenza $16{,}7+95+22+21{,}3=155\unita{ms}$ (\S\ref{sec:lim-latenza});
\item la soglia d'inclinazione oltre la quale un ostacolo alto $a$ scende sotto la quota $y_s=0{,}07\unita{m}$, $\delta\theta=(a-y_s)/\rho$: $0{,}67^\circ$ a $2{,}55\unita{m}$ e $0{,}57^\circ$ a $3\unita{m}$ per $a=0{,}10\unita{m}$;
\item la probabilità di un falso «mobile» per deriva di posa (\S\ref{sec:lim-moto}), controllata con un Monte Carlo indipendente di $2\cdot10^6$ valutazioni: $0{,}0999$ e $0{,}0100$ alle due soglie.
\end{enumerate}
I conti (b)--(e) richiedono solo le quattro operazioni. Le simulazioni non si rifanno dal solo testo: i capitoli~\ref{sec:map}--\ref{sec:occ} e l'appendice~\ref{sec:par} danno regole e parametri, ma non la geometria completa delle disposizioni $G_1$--$G_3$, i parametri del generatore di camminate e dei sette comportamenti, le dieci scene della regola del «libero», il modello del simulatore da scheda (pixel misti, pixel senza dato), l'elenco delle registrazioni pubbliche usate e i semi: per rifarle servono i programmi e i dati del progetto.

\paragraph{Prove a lunga portata.} Le prove del \S\ref{sec:lp} sono programmi in Python con una libreria numerica, sullo stesso calcolatore, e usano semi fissi. Il prototipo simulato si esegue con un solo comando per scena, e tutte le prove e le scene, a corto raggio e a lunga portata, si rifanno insieme:
\begin{quote}\small\ttfamily
sh movi.sh esegui lp\_A1 --config scene/lp\_A1.json --uscita uscite/lp\_A1 --celle\\
sh movi.sh tabellone uscite/lp\_A1/registro.jsonl\\
sh movi.sh convalida uscite/lp\_A1/registro.jsonl\\
sh movi.sh mostra uscite/lp\_A1/registro.jsonl --filmato --uscita uscite/lp\_A1/viste\\
sh prova\_tutto.sh
\end{quote}
Il primo comando simula la scena del drone dietro un edificio e scrive il registro; il secondo dà le quote «dentro la zona» e «dentro l'avviso» con l'intervallo di Wilson; il terzo controlla sul registro i principi (nessun avviso tolto senza guardare); il quarto disegna la pianta e la console dell'operatore. Le scene si chiamano per scenario (drone dietro un edificio, rullaggio, drone sotto la copertura, sensori a bordo di un drone, incrocio, stazione) e per evento (collisione, avaria, fumo e nebbia, drone nel fumo, passaggio ravvicinato). L'esecuzione di tutte le prove e le scene ha dato esito senza errori in $698\unita{s}$. Durate indicative sul calcolatore indicato: circa 35 minuti per le prove simulate per passaggio di tutti gli scenari; $66\unita{s}$ per le prove su traiettorie pubbliche; $88\unita{s}$ per il simulatore dei frammenti con $10\,000$ frammenti per caso; $34{,}5\unita{s}$ per le prove del modello di moto e del «libero». Il costo per ciclo nelle prove simulate è al più $11\unita{ms}$ in mediana con la regola di spegnimento semplice, e arriva a $46\unita{ms}$ in mediana e $340\unita{ms}$ al 95º percentile con quella prudente: non sono tempi di una scheda di bordo. I dati pubblici di traiettoria (tabella~\ref{tab:lp-dati}) si scaricano senza credenziali; le schede dei sensori sono materiale dei costruttori e non sono ridistribuite. Le geometrie complete degli scenari, i semi e i programmi sono nel progetto: dal solo testo si rifanno le distanze d'avviso, le portate del radar per sezione, il «libero» con $n_{\mathrm{eff}}=1$ e gli intervalli di Wilson.

\sloppy\paragraph{Prove del volo ad alta velocità.} Le prove del \S\ref{sec:vo} sono programmi in Python con una libreria numerica, sullo stesso calcolatore, con semi fissi; sviluppo e valutazione usano semi disgiunti, e i parametri sono fissati prima della valutazione. Le uscite si scrivono in file nuovi e non sovrascrivono quelle precedenti. I programmi, nella cartella di lavoro delle prove del progetto, sono:
\begin{itemize}
\item \texttt{eventi\_volo.py} e \texttt{scenari\_volo.json}: gli scenari F1--F6 e il simulatore dei frammenti con l'integratore semi-implicito;
\item \texttt{prove\_volo.py} e \texttt{analizza\_volo.py}: il nucleo con l'inviluppo a $210\unita{m/s}$ (variabile \texttt{VOLO\_REGOLA=v3i}) e il conteggio per passaggio con gli intervalli di Wilson;
\item \texttt{prove\_volo\_v4.py} e \texttt{analizza\_volo\_v4.py}: lo stesso nucleo con la regola di contatto in prova; con \texttt{VOLO\_REGOLA=v3i} si ritrova il nucleo adottato; \texttt{e1\_stabile.py}: la collisione fra multirotori con l'integratore stabile;
\item \texttt{stima\_v4.py} e \texttt{prove\_reali.py}: la confidenza nel tempo e la sua verifica su traiettorie pubbliche; \texttt{stima\_v5.py} e \texttt{prove\_v5.py}: la taratura, con la divisione per mezzo;
\item \texttt{scorta\_f7.py}, \texttt{scorta\_f7r.py} e \texttt{curva\_f7.py}: lo scenario «scorta», la variante con il radar accecato e la curva della confidenza del drone seguito;
\item \texttt{unicita\_v2.py} e \texttt{esempio\_confidenza.py}: il controllo a campione dell'unicità del tempo di soglia su $40\,000$ casi e l'esempio calcolato del \S\ref{sec:vo-tc};
\item \texttt{prototipo/prove/convalida\_volo.py}: il controllo dei principi sui registri delle scene di volo (nessun «libero» sul non osservato, nessun avviso tolto senza guardare, nessun sensore perso che torni affidabile da solo).
\end{itemize}
I comandi principali, con un tempo massimo esplicito e al più 8 processi:
\begin{quote}\small\ttfamily
VOLO\_REGOLA=v3i VOLO\_USCITA=passaggi\_volo\_v3i.jsonl \\\quad timeout 1500 python -I -B prove\_volo.py 60 200 7\\
python -I -B analizza\_volo.py passaggi\_volo\_v3i.jsonl risultati\_volo\_v3i.json\\
timeout 1500 python -I -B prove\_volo\_v4.py 60 200 7\\
python -I -B e1\_stabile.py <modo> 60\\
python -I -B prove\_v5.py <classe> <uscita>\\
python -I -B scorta\_f7.py valutazione 8\\
python -I -B scorta\_f7r.py valutazione 6\\
VOLO\_REGOLA=v3i timeout 600 python -I -B prove\_volo.py --registro prototipo/uscite/video\_volo
\end{quote}
Il primo comando rifà i sei scenari con il nucleo adottato (60 passaggi per cella, 200 per il controllo, 7 processi di lavoro); il secondo conta i passaggi; il terzo rifà le stesse celle con la regola di contatto in prova; gli altri rifanno la collisione fra multirotori, la taratura per classe di mezzo, lo scenario «scorta», la sua variante con il radar accecato e i registri delle scene. Il tempo di calcolo è di 1--8 secondi per passaggio, su CPU, in Python: non è una misura di latenza. Dal solo testo si rifanno la varianza \eqref{eq:vo-sigma}, la confidenza \eqref{eq:vo-c2}, i quantili $k_{0,9}=2{,}146$ (disco) e $2{,}500$ (sfera), il tempo di soglia dell'esempio calcolato, il cammino massimo $L(\tau)$, le distanze d'avviso $d=v\,(L+t_{\mathrm{risp}}+t_{\mathrm{mezzo}})$ e gli intervalli di Wilson.
